\begin{center}
    \textbf{ABSTRACT}
\end{center}

\noindent In partnership with Riachuelo, the team developed the concept of an anti-theft tag that customers release themselves after paying through the store app, with the goal of reducing checkout queues in fashion retail. The market definition set the target segment as mass-market fashion chains that already use self-checkout, and the expansion of Riachuelo's concept stores planned for 2026 was identified as the short-term obtainable market. Consumer research combined a survey with 142 responses and eleven semi-structured interviews. Exploratory clustering of the data identified four groups: service-dependent (39\%), digital self-reliant (21\%), availability-frustrated (9\%) and balanced pragmatists (31\%). Only 24\% of respondents endorsed a purchase process without a manual check at the exit, with full acceptance among the digital self-reliant and near-unanimous rejection among the service-dependent. The findings were translated into eight needs, seven wants and five demands, and seven of these customer voices were deployed into technical requirements through QFD, in which payment-conditioned release, payment time and the clarity of the app flow accounted for about 65\% of the technical importance. Finally, the team drew a freehand sketch of a 98 by 82 by 26 mm device, with a ball-clutch retention mechanism and release authorized by the customer's phone. The results indicate that the solution tends to create value when it preserves the option of human assistance.

\vspace{1em}
\noindent\textbf{Keywords:} fashion retail; anti-theft tag; self-checkout.
